%%% -*- coding: utf-8 -*-
% !TeX root = main-slides.tex

\lecture{Unity Features}{unity-features}
\section{Unity Features}

\showtitlepage{Physics, Materials, Rendering}

%%% Every example is taken from SEE, so the sources named in the listings can
%%% be opened next to the slide.

\subsection{Physics}

\begin{frame}[fragile]{Colliders: the Shape Physics Sees}
  \framelabel{slide:colliders}
  \begin{lstlisting}[basicstyle=\scriptsize\ttfamily]
BoxCollider box = obj.AddComponent<BoxCollider>();   // cheap: six planes
MeshCollider hull = obj.AddComponent<MeshCollider>();// exact, but costly

// Assets/SEE/Game/Layers.cs -- every query is filtered by layer
gameObject.layer = LayerMask.NameToLayer("InteractableGraphObjects");
LayerMask mask = Layers.GraphObjectsLayerMask;
  \end{lstlisting}
  \vspace{1ex}

  \begin{itemize}
    \item A \texttt{Renderer} decides what you see, a \texttt{Collider} what
      physics can touch, and the two are independent. A collider without a
      renderer is invisible but still in the way; a renderer without a
      collider cannot be clicked.
    \item SEE uses mostly \texttt{MeshCollider}, and \texttt{BoxCollider}
      wherever a box will do --- if a node is rendered as a cuboid, a box is both cheaper
      and exact.
    \item Layers are the filter. SEE names five of them in
      \texttt{Layers.cs} --- graph objects and auxiliary objects, each
      interactable or not --- and exposes the matching \texttt{LayerMask}s.
  \end{itemize}
\end{frame}

\begin{frame}[fragile]{Raycasting: What Is Under the Pointer?}
  \framelabel{slide:raycasting}
  \begin{lstlisting}[basicstyle=\scriptsize\ttfamily]
// Assets/SEE/Utils/Raycasting.cs
if (!IsMouseOverGUI()
    && Physics.Raycast(UserPointsTo(), out RaycastHit hit,
                       maxDistance, Layers.GraphObjectsLayerMask))
{
    ...                          // hit.collider, hit.point, hit.distance
}

// all hits along the ray, into a buffer that is kept between frames
numberOfHits = Physics.RaycastNonAlloc(UserPointsTo(), hits,
                                       maxDistance, mask);
  \end{lstlisting}
  \vspace{1ex}

  \begin{itemize}
    \item A ray is an origin and a direction. \texttt{UserPointsTo()} hides
      whether that comes from the mouse or from an XR controller, so the same
      selection code serves the desktop and the headset.
    \item \texttt{RaycastAll} allocates a fresh array per call;
      \texttt{RaycastNonAlloc} fills a buffer SEE keeps, so its loop is bounded
      by \texttt{numberOfHits}.
    \item Without the \texttt{IsMouseOverGUI()} guard a click on a menu would
      also select the city behind it.
  \end{itemize}
\end{frame}

\begin{exerciseframe}{Click to Select}
  Build the smallest possible version of SEE's node selection.
  \vspace{1ex}

  \begin{enumerate}
    \item A scene with a few cubes and a plane. On a left click, raycast from
      the camera through the pointer and log the name of whatever was hit.
    \item Put the cubes on a layer of their own and restrict the query to it,
      so that clicking the plane hits nothing.
    \item Two cubes overlap and you want the nearer one. What does
      \texttt{Physics.Raycast} give you, and what would you need instead to
      choose among all hits, as \texttt{RaycastLowestNode} does?
  \end{enumerate}
\end{exerciseframe}

\begin{solutionframe}{Click to Select}
  \begin{lstlisting}[basicstyle=\scriptsize\ttfamily]
private void Update()
{
    if (Input.GetMouseButtonDown(0)
        && Physics.Raycast(Camera.main.ScreenPointToRay(Input.mousePosition),
                           out RaycastHit hit, 100.0f, cubeMask))
    {
        Debug.Log($"{hit.collider.name} at {hit.point}\n", hit.collider);
    }
}
  \end{lstlisting}
  \vspace{0.5ex}

  \begin{itemize}
    \item \texttt{cubeMask} is a \texttt{LayerMask} field set in the inspector,
      or \texttt{1 << LayerMask.NameToLayer("Nodes")} in code. Passing no mask
      means \emph{everything}, the plane included.
    \item \texttt{Physics.Raycast} already returns the \emph{nearest} hit, so
      part three needs nothing extra. Choosing a different one --- the deepest
      node rather than the closest surface --- is what forces
      \texttt{RaycastNonAlloc} and a loop over all hits.
  \end{itemize}
\end{solutionframe}

\subsection{Materials and Shaders}

\begin{frame}[fragile]{A Material Is a Shader Plus Values}
  \begin{lstlisting}[basicstyle=\scriptsize\ttfamily]
Material m = new(Shader.Find("Universal Render Pipeline/Lit"));
m.color = Color.red;              // a property of that shader
m.SetFloat("_Metallic", 0.5f);    // by name: a string lookup every call

// Assets/SEE/Game/Portal.cs -- the name is looked up once
private const string portalPropertyName = "_Portal";
private static readonly int portalProp =
    Shader.PropertyToID(portalPropertyName);
  \end{lstlisting}
  \vspace{1ex}

  \begin{itemize}
    \item The \emph{shader} is the program that runs on the GPU; the
      \emph{material} is an instance of it with values filled in. Two red
      cubes and two blue ones need one shader and two materials.
    \item Shader properties are addressed by name, and Unity turns that name
      into an integer. \texttt{Shader.PropertyToID} does it once into a
      \texttt{static readonly int} --- the idiom SEE uses for
      \texttt{\_Portal}, its own property that clips a city to its table.
    \item Which shaders exist depends on the render pipeline:
      \texttt{Universal Render Pipeline/Lit}, \texttt{Simple Lit} and
      \texttt{Unlit} are URP's, and SEE adds its own under
      \texttt{Shader Graphs/}.
  \end{itemize}
\end{frame}

\begin{frame}[fragile]{Sharing Beats Instancing}
  \begin{lstlisting}[basicstyle=\scriptsize\ttfamily]
// Assets/SEE/GameObjects/Factories/MaterialsFactory.cs
public void SetSharedMaterial(Renderer renderer, int index)
{
    renderer.sharedMaterial = Get(Mathf.Clamp(index, 0, ...));
}

renderer.material.color = Color.red;        // clones the material, silently
renderer.sharedMaterial = existingMaterial; // assigns the shared asset
  \end{lstlisting}
  \vspace{1ex}

  \begin{itemize}
    \item Reading \texttt{renderer.material} does not read: it \emph{clones}
      on first access and hands you a private copy. One per object, no longer
      batched together, and never freed by itself.
    \item That is why \texttt{MaterialsFactory} builds one material per color
      of a \texttt{ColorRange} up front and then only hands out indices. A
      city of thousands of blocks draws with a handful of materials.
    \item Editing a \texttt{sharedMaterial} in play mode edits the asset on
      disk. It is the same trap as a \texttt{ScriptableObject}, and for the
      same reason.
  \end{itemize}
\end{frame}

\begin{exerciseframe}{One Red Cube, Then a Thousand}
  \begin{enumerate}
    \item Ten cubes share one material. Turn exactly one of them red at run
      time. Then inspect the other nine and the material asset --- what
      happened, and what would have happened had you written
      \texttt{sharedMaterial} instead?
    \item Now color a thousand cubes along a gradient from green to red,
      without creating a thousand materials. How many do you need, and what
      does each cube have to store?
  \end{enumerate}
  \vspace{1ex}

  \hint{Part two is what a code city does for every metric it maps to color.}
\end{exerciseframe}

\begin{solutionframe}{One Red Cube, Then a Thousand}
  \begin{itemize}
    \item \texttt{renderer.material.color = Color.red} clones the material for
      that one cube: the other nine keep the original and the asset on disk is
      untouched. \texttt{renderer.sharedMaterial.color = Color.red} would have
      turned all ten red --- and saved the change into the asset.
    \item You need as many materials as you want distinguishable steps, not as
      many as there are cubes. Build the gradient once, keep the materials in
      an array, and give every cube an index into it; that is exactly
      \texttt{MaterialsFactory} with its \texttt{ColorRange}. Each cube stores
      an integer, and all cubes sharing a step still batch.
  \end{itemize}
  \vspace{1ex}

  \hint{The alternative is one material with per-instance data through a
  \texttt{MaterialPropertyBlock}. It keeps a single material but gives up
  static batching, so measure before choosing.}
\end{solutionframe}

\subsection{The Render Pipeline}

\begin{frame}{The Pipeline Is an Asset, Too}
  \begin{itemize}
    \item URP --- the \emph{Universal Render Pipeline} --- replaces Unity's
      built-in renderer with one you configure rather than compile. The
      configuration is a \texttt{ScriptableObject}, so everything from
      \slideref{slide:scriptableobject} applies.
    \item SEE ships one pipeline asset per quality level, each with its own
      renderer: \texttt{Fastest}, \texttt{Fast}, \texttt{Simple},
      \texttt{Good}, \texttt{Beautiful}, \texttt{Fantastic}. Switching quality
      switches the whole renderer, not just a few numbers.
    \item Post-processing is a second asset: a \texttt{Volume} in the scene
      points at a profile, \texttt{DefaultVolumeProfile} or SEE's
      \texttt{BeautifyProfile}, and the effects come from there.
  \end{itemize}
  \vspace{2ex}

  \hint{A shader written for one pipeline does not work in another. An object
  that turns magenta is almost always a material whose shader belongs to the
  built-in pipeline, not to URP.}
\end{frame}

\begin{exerciseframe}{Magenta and Other Mysteries}
  \begin{enumerate}
    \item You import a free asset from the store and its objects come out
      magenta. What happened, and what are the two ways to fix it?
    \item Shadows should reach further in the VR build than on the desktop,
      and nowhere else should change. Which file do you edit, and which do you
      leave alone?
    \item \texttt{Assets/Settings} holds a pipeline asset \emph{and} a forward
      renderer per quality level. Why two files rather than one?
  \end{enumerate}
\end{exerciseframe}

\begin{solutionframe}{Magenta and Other Mysteries}
  \begin{enumerate}
    \item Magenta is Unity's error shader: the material asks for a shader that
      the active pipeline does not provide, typically the built-in
      \texttt{Standard}. Either upgrade the materials --- Unity has a
      converter for it --- or replace the shader by
      \texttt{Universal Render Pipeline/Lit} by hand.
    \item The shadow distance lives in the pipeline asset of the quality level
      the VR build runs at, so edit that one asset. The other five levels and
      the scene stay untouched --- which is the point of having one asset per
      level.
    \item Because they answer different questions. The pipeline asset says
      \emph{how much}: shadow distance, HDR, MSAA, render scale. The renderer
      says \emph{how}: the passes and renderer features that produce a frame.
      Several pipeline assets can point at the same renderer.
  \end{enumerate}
\end{solutionframe}

%%% Local Variables:
%%% mode: latex
%%% TeX-master: "main-slides"
%%% mode: reftex
%%% mode: flyspell
%%% End:
